\clearpage
\fancyhead[L]{\small CONNOISSANCE DES TEMS.}
\hypertarget{native-522-lower}{}\label{native-522-lower}
\begin{center}\small Göttingische gelehrte Anzeigen. Stück 95. S. 939..943. 17. Juni 1809.\\[3pt]\rule{.72\linewidth}{.3pt}\end{center}

\emph{Connoissance des tems ou des mouvemens célestes, à l'usage des astronomes et des navigateurs, pour l'an} 1810; \emph{publiée par le bureau des longitudes. De l'imprimerie impériale} 1808. Großoktav. Paris.

Wir übergehen den Kalender und die übrigen ständigen Artikel, deren Einrichtung unverändert geblieben ist, mit Stillschweigen. Wir bemerken nur, dass die Tafel der geopraphischen Längen und Breiten auch diesmal einige nicht unbedeutende Ergänzungen erhalten hat, vor allem \SourceLeaf{523}mehrere wertvolle neue Bestimmungen von der Insel Zypern, aus Arabien und vom Roten Meer. Die \emph{Additions}, die der eigentliche Anlass für diese Anzeige sind, beginnen wie in den beiden vorigen Jahrgängen mit den Beobachtungen, die Herr \textsc{Bouvard} im Jahr 1806 auf der kaiserlichen Sternwarte angestellt hat. Sie füllen 80 Seiten. Außer Sonne und Mond finden wir darunter auch sämtliche älteren Planeten, diesmal aber gar keine Beobachtungen von Ceres, Pallas oder Juno. Den von \textsc{Pons} im November in Marseille entdeckten Kometen beobachtete \textsc{Bouvard} vom 21. November bis zum 19. Dezember. Verfinsterungen von Jupitermonden sind ziemlich zahlreich. Sternbedeckungen gibt es insgesamt fünf, und zwar nur Eintritte ($g$ im Löwen am 9. Januar, $\zeta$ im Krebs am 1. März, $\theta$ im Ophiuchus am 1. Juni, $\xi$ in den Zwillingen am 8. September, 19 in den Fischen am 20. November). — Darauf folgen chinesische Beobachtungen seit dem Jahr 147 vor unserer Zeitrechnung, die Pater \textsc{Gaubil} im Jahr 1749 eingesandt hat. Die Beobachtungen beschränken sich auf enge Konjunktionen und Bedeckungen des Mondes und der Planeten untereinander und mit Fixsternen. Nähere Umstände fehlen, und meistens ist nur der Tag ohne die Stunde angegeben. Viel praktischer Nutzen wird sich daraus also kaum ziehen lassen. — Eine literarische Notiz von \textsc{J. C. Burckhardt} über die Verdienste von \textsc{Hevel} und \textsc{Dörfel} um die Theorie der Kometenbewegung. — Sechste und letzte Sammlung der Beobachtungen von \textsc{Messier}, von 1752 bis Ende 1759. Sie enthält zwei Mondfinsternisse, viele Sternbedeckungen durch den Mond, noch viel zahlreichere Verfinsterungen von Jupitermonden und zwei Kometen. Einer davon, der \textsc{Halley}sche Komet, wurde am 21. Januar 1759 zuerst aufgefunden. Auf Verlangen seines Lehrers \textsc{Delisle} musste \textsc{Messier} die Entdeckung bis Anfang April geheim halten. — Beobachtung der oberen Konjunktion des Merkur und der unteren der Venus im September und Oktober 1807 durch \textsc{Vidal} in Mirepoix. — Beobachtung des großen Kometen von 1807 durch denselben Beobachter vom 27. November bis zum 4. März 1808. — Beobachtungen desselben Kometen an der königlichen Sternwarte in Lissabon durch \textsc{Paul Ciera} (vom 7. Oktober bis zum 29. November) und in Bremen durch \textsc{Olbers} (vom 8. Oktober bis zum 14. Februar 1808). — Messung eines Erdmeridianbogens und eines auf dem Meridian senkrecht stehenden Grades in Ostindien durch Brigademajor \textsc{William Lambton}. Es handelt sich um einen Auszug aus den Memoirs von Calcutta. Die Basis maß 40006 englische Fuß. Die Endpunkte des gemessenen Meridianbogens waren

\begin{center}\begin{tabular}{@{}lll@{}}\emph{Paudree} & Breite & $13^\circ\,19'\,49''018$ \\ \emph{Trivandepcrum} & Breite & $11\quad44\quad52.\,59$ \\ \end{tabular}\end{center}
\noindent Daraus leitete man für den Breitengrad 60495 Fathoms oder 56763 Toisen ab. Für den Längengrad fand man 61061 Fathoms oder 57294 Toisen. Wie groß der Bogen war, aus dem der Längengrad bestimmt wurde, ist nicht angegeben. Daher können wir nicht beurteilen, welchen Grad an Genauigkeit man ihm zuschreiben darf. Der Vergleich der beiden würde die zu starke Abplattung $\frac{1}{205,67}$ ergeben. Mit der französischen Gradmessung verglichen, würde der Breitengrad die Abplattung $\frac{1}{817}$ ergeben. — Beobachtungen an der Lissaboner Sternwarte durch \textsc{Paul Ciera} (außer drei Bedeckungen des Sterns $\alpha^2$ im Krebs durch den Mond, 1807 am 20. Februar, am 16. April und am 14. Oktober, nur Verfinsterungen von Jupitermonden). — Methode von \textsc{Delambre} zur Berechnung der Korrekturen der Durchgänge am Passagenfernrohr. Der Aufsatz ist sehr weitschweifig, mag aber weniger geübten Praktikern nützen. Herr \textsc{Delambre} scheint es für eine neue Bemerkung zu halten, dass man das Passagenfernrohr allein durch den Vergleich mehrerer Durchgänge mit den Rektaszensionen lediglich in einen Stundenkreis bringen kann, ohne es damit in den Meridian zu bringen, wenn man weder die Libelle noch die absolute Sternzeit noch sonst etwas heranzieht, das sich auf den Zenit des Beobachters bezieht. Wir können aber nicht glauben, dass etwas so Offenkundiges je einem nachdenkenden Astronomen entgangen sein kann. — Die Mondfinsternis vom 4. Januar 1806, beobachtet von \textsc{Vidal} in Mirepoix. — Beobachtungen an der kaiserlichen Sternwarte in Marseille von \textsc{Thulis}, von 1796..\SourceLeaf{524}1806. — Beobachtung des großen Kometen von 1807 in Montauban durch \textsc{Duc-la Chapelle}. — Beobachtungen in Viviers durch \textsc{Flaugergues} im Jahr 1807. — Die allgemeinen Aberrations- und Nutationstafeln des Herrn \textsc{von Zach}, aus seinem bekannten Werk abgedruckt. Ebenso die allgemeinen Aberrations- und Nutationstafeln von \textsc{Gauss}, aus der Monatl. Correspondenz abgedruckt, zusammen mit einer Herleitung der Formeln, auf denen sie beruhen. — Geschichte der Astronomie für 1808. Sie besteht diesmal nur aus Auszügen aus neu erschienenen astronomischen Werken. In \textsc{von Zach}s \emph{Tabulae speciales} nimmt Herr \textsc{Delambre} die Gelegenheit wahr, die Formeln mitzuteilen, auf denen seine in diesem Werk erneut benutzten Aberrationstafeln für die Planeten beruhen. Es folgen \textsc{Cagnoli}s \emph{Sternverzeichnis}, \textsc{Monteiro-da-Rocha}s \emph{Mémoires sur l'astronomie pratique}, aus dem Portugiesischen ins Französische übersetzt, die französische Übersetzung von \textsc{Cagnoli}s Trigonometrie nach der zweiten Ausgabe, die erste Lieferung des astronomischen Teils der Reise von \textsc{Alexander von Humboldt}, \textsc{L. W. Pfaff}s \emph{de tubo Culminatorio Dorpatensi brevis narratio}, \emph{Dorpati} 1808, und die \emph{Efemeridi astronomiche di Milano per l'anno} 1809. An dieser Aufzählung sieht man, dass die \emph{Connoissance des tems} für Astronomen weiterhin ein sehr wertvolles Nachschlagewerk ist, auch wenn dieser Band etwas weniger Originalaufsätze enthält als manche seiner Vorgänger. Man wird darin nicht vergeblich nach Nachrichten über alles besonders Wissenswerte suchen, was in der Astronomie geschieht. — Auf diese Auszüge folgt noch ein Bericht über die Endergebnisse der französischen Gradmessung, die bis zur Insel Formentera fortgesetzt wurde. Dabei wurde ein Dreieck benutzt, dessen eine Seite 82555 Toisen maß. Der ganze Bogen von Dünkirchen bis Formentera ist dadurch auf $13^\circ\,12'\,22''$ angewachsen, und der Meter verringert sich dadurch von 443,296 Linien auf 443,2958 Linien, ändert sich also so gut wie gar nicht. Den Abschluss des Bandes bilden einige Verbesserungen zu den vom Bureau des Longitudes herausgegebenen Sonnen- und Mondtafeln, eine neu berechnete Tafel für die Zeitgleichung für 1800 zusammen mit den säkularen Änderungen und die meteorologischen Beobachtungen, die \textsc{Bouvard} im Jahr 1806 angestellt hat.
